\begin{abstract}
Agents operating over evolving time-series data repeatedly face the same analytical structure under new release periods. TimeSage-EV shows that current agents still exhibit recurring temporal-validity failures and recurring exogenous-context failures in this setting. The benchmark's self-evolving TimeSage-1.0 harness adds a reusable skill library and improves aggregate performance while reducing token cost to 0.82$\times$ the baseline. We study the mechanism behind these gains. Our hypothesis is that a large class of recurring failures is caused by missing reusable procedures for cutoff-aware temporal reasoning and exogenous-context integration. We formulate targeted temporal skills as a causal intervention. The intervention injects a compact procedure for a known failure mode while keeping the future endpoint and evidence fixed. A matched generic skill controls for the benefit of receiving any reusable procedure. A no-skill condition provides the baseline. This design distinguishes targeted mechanism repair from generic retrieval assistance. The source benchmark exposes an unusually clean intervention surface because accepted skills become callable tools starting from the next period. We specify a skill contract for temporal cutoff enforcement. A second contract handles release-to-period alignment. A third contract grounds exogenous events. The resulting protocol predicts that targeted skills will remove matched recurring failures more strongly than generic skills while preserving chronological evaluation. This work advances self-evolving time-series agents from aggregate skill ablation toward mechanism-specific causal evaluation. It also provides a general recipe for testing whether recurring agent failures arise from missing reusable procedures.
\end{abstract}
